\documentclass[11pt,a4paper]{article}
\usepackage[T1]{fontenc}
\usepackage{lmodern}
\usepackage[margin=26mm]{geometry}
\usepackage{amsmath,amssymb,booktabs,tabularx,array}
\usepackage{microtype}
\usepackage{xcolor}
\usepackage{fancyhdr}
\usepackage[hidelinks]{hyperref}
\definecolor{accent}{RGB}{25,67,104}
\hypersetup{colorlinks=true,linkcolor=accent,urlcolor=accent,citecolor=accent}
\pagestyle{fancy}
\fancyhf{}
\lhead{\texttt{net-center} | Technical note}
\rhead{Network location}
\cfoot{\thepage}
\setlength{\headheight}{15pt}
\setlength{\parindent}{0pt}
\setlength{\parskip}{5pt}
\newcommand{\code}[1]{\texttt{\detokenize{#1}}}
\newcommand{\dist}{d_G}
\newcommand{\minsite}{\min_{s\in S}}
\title{\textbf{Network Location Optimization and\\Adaptive Spatial Representation}\\[8pt]\large Technical note for \texttt{net-center}}
\author{James Matthew Miraflor\\\small Scientific Computing Laboratory, Department of Computer Science\\\small University of the Philippines Diliman}
\date{26 November 2026\\\small Implementation baseline: \texttt{e7710b9b436f1} (unreleased adaptive expansion)}
\begin{document}
\maketitle
\begin{abstract}
\noindent This note describes the original exact weighted 1-median, vertex 1-center and absolute 1-center methods together with the adaptive location workflow. The latter supports finite demand attached to edge interiors, independent groups, finite-candidate multiple medians, explicit component allocation and criterion-based representative selection. Exactness is restricted to the modeled problem and stated numerical tolerances. Multi-site BUILD/SWAP is a heuristic; bounded exhaustive enumeration is available. The algorithms are established methods and the selection criteria and reference support are explicit implementation choices, not claims of new mathematical theorems.
\end{abstract}
\setcounter{tocdepth}{1}
\tableofcontents
\clearpage
\section{Model and scope}
Let $G=(V,E)$ be a finite undirected metric graph with positive edge lengths $L_e$. The GIS interface takes a line layer in a projected metre CRS; numeric callers supply a consistent positive symmetric routing metric. Let $q_i\in V$ be the attached location of finite observation $i$, $w_i\ge0$ its weight, and $a_i\ge0$ its fixed off-network connector length. Let $S$ be an eligible set of sites, and $I$ the set of modeled rows. Demand groups are independent optimization problems, even when sharing a graph.

\begin{tabularx}{\linewidth}{@{}p{.27\linewidth} X X@{}}
\toprule
Problem & Candidate domain & Method / guarantee \\
\midrule
Weighted 1-median & Finite vertices & Exact distance reduction \\
Vertex 1-center & Finite vertices & Exact maximum-distance reduction \\
Absolute 1-center & Physical metric edges & Exact piecewise envelope sweep, numerical tolerance \\
Weighted multiple medians & Finite vertices, demand or restricted set & BUILD/SWAP heuristic or bounded exact enumeration \\
Component allocation & Evaluated cost curves & Exact dynamic program over those curves \\
Adaptive site counts & Evaluated solutions & Criterion-based choice, conditional on curves \\
\bottomrule
\end{tabularx}

Define the distance to the nearest eligible selected facility as $r_i(S)=\minsite \dist(q_i,s)$ and access as $b_i(S)=a_i+r_i(S)$. For positive modeled mass $W=\sum_{i\in I}w_i$,
\begin{align*}
J(S)&=\sum_{i\in I}w_i r_i(S), &
J_a(S)&=J(S)+\sum_{i\in I} w_i a_i,\\
R(S)&=\max_{i\in I}r_i(S), & R_a(S)&=\max_{i\in I}b_i(S).
\end{align*}
For fixed modeled demand and attachment, the connector term of $J_a$ is constant in $S$. The minimax optimum can move when offsets are included. Weights affect median costs, not minimax radii. Zero-weight rows still require coverage and enter maxima. Site budgets and grouping are explicit; the continuous minimax wrapper is not a general absolute $k$-center algorithm.

The original CLI and \code{solve} retain the legacy single-site model. The adaptive CLI and its submodule APIs are additional interfaces. The term ``network center'' is not an arithmetic centroid: Euclidean means minimize squared Euclidean distances, a different objective. A network squared-distance Fr\'echet mean is not implemented.
\section{Topology, attachment and the metric}
Input line geometry is not itself a routing graph. Source-shared vertices recover junctions, including endpoint/interior and self-touching cases. Geometric crossings without shared vertices do not connect by default. Optional planar noding assumes all crossings are junctions and is inappropriate for grade separation. The graph builder does not infer one-way, access, bridge, tunnel or turn attributes. A transitive coordinate tolerance identifies coincident endpoints; rings are split, slivers removed and node IDs compacted. The physical edge list retains distinct parallel geometries; CSR stores the minimum length for each endpoint pair.

Legacy attachment uses nearest existing nodes via a k-d tree. Adaptive attachment projects each observation to the nearest physical LineString (deterministic lowest edge ID for exact ties). If an accepted offset $t\in(0,L_e)$ is new, split the edge there, preserving original order, lineage and total length. Duplicate offsets are inserted only once. Original shortest paths remain unchanged: any path along an old edge can be expanded into consecutive segments of equal total length, and consecutive subsegments contract to the old edge without changing the metric. New inserted vertices are genuine locations on the physical network. Different crossing edges remain distinct. Finite attached points do not constitute continuous demand density along the entire graph.

Ordered endpoint correspondence, physical lengths, minimum-parallel-edge CSR and retained transitive endpoint chains are validated. Stale snap records are rejected. Node representatives can differ from physical endpoints within the original topology tolerance. No new intercomponent links are invented. Coordinates must use the declared metre CRS; for extensive study regions choose a suitable projection rather than assuming one UTM zone is adequate.

\textbf{Component retention before versus after snapping.} The legacy workflow retains the largest component by physical length first. Adaptive GIS retains all components unless \code{--keep-largest-component} is selected. Pre-filtering can move an observation farther from its source coordinate. The adaptive \code{largest} policy instead snaps first, then excludes rows outside the largest component. These are different demand models. Both must respect maximum snap distance; adaptive exclusion is explicit and audited. The full graph and original physical lengths determine GIS ranking, not the number of inserted vertices. CSR-only calls default to the collapsed routing weights; longer parallel links cannot be inferred from the CSR.
\section{Original single-center algorithms}
\subsection{Distance matrix and vertex solutions}
Dijkstra's algorithm \cite{dijkstra} supplies $D_{iv}=\dist(q_i,v)$. Repeated demand nodes are coalesced with summed weights for the median; duplicates do not change an unweighted maximum. The legacy solver builds the demand-by-vertex matrix, then computes
\[
 v_M\in\arg\min_v\sum_iw_iD_{iv},\qquad v_C\in\arg\min_v\max_iD_{iv}.
\]
Hakimi's vertex-optimality result permits a median optimum at a network vertex for finite vertex demand \cite{hakimi64,hakimi65}. Restricted candidates change the feasible problem. For minimax a continuous optimum may be interior to an edge.
\subsection{Absolute 1-center: edgewise minimization}
For edge $e=(u,v)$ of length $L$, the distance to a location at offset $t$ from $u$ is
\[
 h_i(t)=\min\{D_{iu}+t,D_{iv}+L-t\},\quad 0\le t\le L.
\]
The switch point is $t_i=(D_{iv}-D_{iu}+L)/2$. Shortest-path triangle inequalities give $0\le t_i\le L$ (the implementation clips numerical excursions). The objective is $F(t)=\max_i h_i(t)$. Sort switch points $t_{(i)}$. On an interval between adjacent switches, each demand uses a fixed endpoint route. The envelope becomes
\[
 F(t)=\max\{A+t,B+L-t\},
\]
where $A$ is the maximum $D_{iu}$ among unswitched demands and $B$ the maximum $D_{iv}$ among switched demands. Empty maxima are $-\infty$. The interval minimum is found at the clipped intersection $(B+L-A)/2$, with endpoint cases when one branch is absent. Prefix and suffix running maxima evaluate all intervals after sorting. Compare the best interior location with the best vertex. This is the classical edgewise construction \cite{kariv,handler} implemented using vectorized arrays.

For the grouped absolute-center wrapper, when connector costs are included, replace both endpoint distances for row $i$ by $D_{iu}+a_i$ and $D_{iv}+a_i$. The switch difference cancels $a_i$, so the same sweep applies. At duplicate nodes use the largest offset for the maximum; other rows still receive individual assignment/access results. The final routed radius is recomputed and checked.
\subsection{Safe lower bounds}
Let $E(u)=\max_i D_{iu}$, or the corresponding connector-inclusive maximum. The current best vertex gives an incumbent radius $R_0$. Any point of edge $e$ satisfies $F(t)\ge E(u)-t$ and $F(t)\ge E(v)-(L-t)$. Thus
\[
 \mathrm{LB}_1(e)=\tfrac12(E(u)+E(v)-L)\le\min_t F(t).
\]
A second bound is $\mathrm{LB}_2(e)=\max_i\min(D_{iu},D_{iv})$, using offset-adjusted distances when appropriate. It is applied to survivors of the first bound. Both can discard edges unable to improve the incumbent; numerical margins protect against floating-point pruning errors. The solver is exact up to storage precision and its comparison tolerance, not exact real arithmetic. The original derivation and historical optimizations remain in \code{docs/ALGORITHMS.md}.
\section{Adaptive single-site and finite-candidate optimization}
\subsection{Streaming median and vertex minimax}
The adaptive finite-site solver avoids a persistent demand-by-node matrix for single-site subproblems. Accumulate a weighted sum or row-wise maximum over bounded Dijkstra batches on each component, then choose the best eligible vertex. Zero-weight demand rows are omitted from the weighted sum to avoid undefined products with infinity; they remain in assignments and minimax. Temporary memory is a target, with at least one complete row required.

For a median with positive mass $W$ and a feasible reference candidate $p$, the triangle inequality gives
\[
 J(v)\ge W\dist(p,v)-J(p).
\]
Indeed $\dist(q_i,v)\ge\dist(p,v)-\dist(p,q_i)$ and summation yields the bound. An improving vertex cannot be farther than $2J(p)/W$ from $p$. For candidates within that radius, every positive-weight demand distance is bounded above by $\max_i\dist(p,q_i)+2J(p)/W$. This supplies a safe Dijkstra search cutoff, with a numerical margin. It is an engineering acceleration, not a change to the objective. The optional bound is not a multi-site pruning theorem.
\subsection{Candidate universes and groups}
The \code{vertices} domain is all supplied vertices in a numerical call, or original vertices plus the particular group's own accepted demand vertices in the GIS workflow. The \code{demand} domain restricts to distinct demand positions (medoids). The \code{restricted} domain uses explicit candidate vertices or a snapped candidate point layer. A shared augmented graph can include positions inserted by unrelated groups for routing, but those positions are not automatically facility candidates. Exact and heuristic multi-site algorithms operate on the same finite set. Group budgets apply independently.
\subsection{Multi-site cost curves}
For each component $c$ and feasible count $j$, compute a cost $J_c(j)$ using an eligible candidate set. In the heuristic, BUILD greedily adds the site with maximum improvement in weighted nearest-site distance (including the best single site as its first step). SWAP examines replacements and executes a best improving exchange for at most the specified number of passes. This is the family of interchange/PAM-style heuristics \cite{teitz,kaufman}; global optimality is not guaranteed, Each larger curve point extends the previous selected set before improving swaps. Deterministic candidate order and strict improvement tolerance control ties.

Exact enumeration evaluates each eligible $j$-subset. If the sum of binomial counts for the requested curve exceeds the configured per-component cap, fail before allocating its dense matrix. This is exact only over finite candidates and evaluated site counts, not over arbitrary continuous facility positions. Both multi-site engines retain a dense unique-demand by candidate matrix; the matrix memory gate does not bound total process memory.
\section{Disconnected components and allocation}
For a group, let $C$ denote the number of modeled demand-bearing components. Each requires a site for finite coverage; no travel is assumed between components. The policies are:
\begin{description}
\item[Strict] Require one populated component. Other components may remain in the graph.
\item[Largest] Exclude demand outside the largest component after snapping; require remaining positive mass.
\item[Per component] Solve one site independently in each, within the group cap.
\item[Allocate] Distribute the group budget over populated components, with at least one site per component.
\end{description}
In a median group with positive total modeled mass, a component with only zero-weight rows receives exactly one site minimizing its largest connector-inclusive access. Its median cost is zero. This is a secondary coverage convention, not a consequence of median optimality; it can make an otherwise feasible fixed-count problem infeasible. Minimax rows are unweighted and may include zero-weight observations. All-zero modeled groups are rejected.

Define the dynamic programming state for the first $c$ components and total $K$ sites by
\[
 F_c(K)=\min_{j\ge1}\bigl\{F_{c-1}(K-j)+J_c(j)\bigr\},
\]
with $F_0(0)=0$ and other initial states infeasible. Only available curve points and counts up to the group cap enter the recurrence. The dynamic program retains the selected site tuple at each state. If all component curves are exact, so is the fixed-$K$ solution under the modeled constraints. If they are heuristic, this only optimizes allocation among the computed sets. Minimax per-component aggregation instead takes the maximum of independently optimized radii. General multi-site minimax allocation is not implemented.
\section{Site-count selection}
For median objectives, let $J(K)=F_C(K)$ be the evaluated network-cost frontier, $n$ the number of distinct modeled demand vertices, and $W$ total modeled mass. Adaptive selections require \code{k=None} in the Python API (omit \code{--k} in CLI). Minimax supports fixed single-site subproblems only.
\subsection{Fixed, balanced, budget and penalized}
\textbf{Fixed:} choose the requested feasible $K$; with \code{per_component}, use exactly $C$ sites. \textbf{Balanced:} choose an evaluated $C\le K\le n$ minimizing
\[
 \frac{J(K)}{J(C)}+\lambda\frac{K-C}{\max(1,n-C)},
\]
where the normalized loss is zero when $J(C)=0$, and $\lambda\ge0$ (default 1). The denominator counts distinct observed vertices, not candidate vertices or total rows. Smaller $K$ wins exact score ties. The criterion expresses a modeling tradeoff; it is not an inferential theorem.

\textbf{Access budget:} choose the first evaluated $K$ satisfying
\[
 \frac{J(K)+\sum_i w_i a_i}{W}\le B+10^{-12}\max(1,B).
\]
The routed access is checked again; failure is not a proof that no feasible facility set exists outside the evaluated curves. \textbf{Penalized:} minimize $J(K)+fK$ where $f\ge0$ is a facility cost in weight-metres. The connector term is constant; the network-only selection score is reported separately from the location objective.
\subsection{Network-medoid Gap variant}
The Gap statistic of Tibshirani, Walther and Hastie \cite{gap} compares observed dispersion against a reference distribution. Here $W_j=J_c(j)$ is the weighted sum of \emph{unsquared} network nearest-medoid distances (power 1), not squared Euclidean dispersion. Draw reference observations with replacement from eligible component vertices, optionally inside the observed demand bounding box. GIS uses original pre-augmentation vertices; preserve individual row weights. The null samples vertices, not uniform physical edge length. Fit heuristic medoids to observed and reference samples.

For $B$ reference replicates, set
\begin{align*}
 g_j&=B^{-1}\sum_{b=1}^{B}\log W^*_{jb}-\log W_j,\\
 s_j&=\operatorname{sd}_b(\log W^*_{jb})\sqrt{1+1/B}.
\end{align*}
Logs apply a common floor $\max(\text{float64 tiny},10^{-12}J_c(1))$. Choose the smallest tested $j$ with $g_j\ge g_{j+1}-s_{j+1}$; otherwise report upper-bound contact. Insufficient reference support, degenerate loss or a rejected minimum-member split retains one site with explicit audit. Component selections are independent and a global cap violation fails. This is a network variant, not a reproduction of R \code{clusGap}; see its \code{d.power} discussion \cite{clusgap}. Results depend on topology, vertex density, extent, weights, seed and heuristic. Do not interpret a cap-limited choice as conclusive evidence for a true number of clusters.
\section{Representation and diagnostics}
For selected sites $S$, assign each modeled row to its nearest site, breaking distance ties by increasing site node ID in the finite-site workflow. For site $s$, let $I_s$ be its assigned rows. The represented mass and mean access are
\[
 W_s=\sum_{i\in I_s} w_i,\qquad
 \bar b_s=\frac{\sum_{i\in I_s}w_i b_i}{W_s}\quad(W_s>0).
\]
Report also assigned row counts and maximum access, including modeled zero-weight rows. An empty site has undefined maximum; a zero-mass site has undefined weighted mean. Weighted percentiles are inverse empirical CDF values on positive modeled mass; the denominator for the overall weighted mean is modeled mass, not all input mass. Coverage records the modeled fraction of original mass and rows. Unassigned rows are explicit, not zero-distance observations. The GIS result reports both attached and full-input coverage. JSON converts nonfinite values to null, preserving statuses that identify exclusions.

The adaptive schema is \code{net-center-adaptive/5}. Invocation settings, resolved finite-site tuning options and component execution diagnostics are separate. The latter identify the actual solver, evaluated counts, enabled streaming pruning and generated Gap reference replicates, not a claim that every configured parameter affected execution. The continuous wrapper has separate location and radius diagnostics. A continuous site is a node or a physical edge and offset; edge lineage traces it to the original network. Final routes are recomputed to verify the evaluated objective. All groups must finish before atomic no-clobber JSON publication via hard link. Input data fingerprints must be recorded externally for full reproducibility.
\section{Implementation and computational limits}
\begin{tabularx}{\linewidth}{@{}p{.35\linewidth}X@{}}
\toprule
Module & Responsibility \\
\midrule
\code{graph.py}, \code{topology.py} & CRS and physical topology; minimum parallel edge CSR \\
\code{distances.py}, \code{center.py}, \code{solve.py} & Legacy Dijkstra matrix, exact single-site solvers, orchestration \\
\code{edge_snap.py} & Physical geometry validation, snapping and edge augmentation \\
\code{adaptive.py} & Component cache, finite-site optimization, allocation and selection \\
\code{absolute_groups.py} & Validated component wrapper around absolute-center sweep \\
\code{cli.py}, \code{adaptive_cli.py} & Distinct legacy and adaptive GIS interfaces \\
\bottomrule
\end{tabularx}

With $N$ vertices, $M$ edges and $Q$ unique demand positions, the legacy matrix requires $O(QN)$ stored distance values. A standard heap Dijkstra estimate is $O(M+N\log N)$ per source (actual SciPy runtime depends on its implementation and graph). Absolute sweep worst-case work is $O(MQ\log Q)$ after distances, reduced by lower bounds and blocked evaluation. Legacy float32 storage changes numerical precision; weighted sums accumulate in float64. Joblib parallelism and legacy memory controls do not automatically apply to adaptive multi-site algorithms.

Adaptive single-site streaming avoids storing the full matrix but retains graph/subgraphs, score arrays and at least one distance row. Multi-site candidate count $P$ requires $O(QP)$ persistent distance values and temporary routing blocks. Exact subset enumeration scales combinatorially; BUILD/SWAP work grows with candidates, demand, site counts and swap passes. Gap repeats reference matrix/curve calculations. Augmentation traverses physical edges; no national-scale guarantee is claimed. A block-memory target is not a bound on RSS. General graph contraction, out-of-core matrices, directed/asymmetric costs, capacity constraints, general weighted minimax, continuous demand density and multi-facility minimax are not implemented.
\section{Validation and reproducibility}
The baseline merged commit \code{e7710b9b436f1b344d60d2327babb42e0e36f266} passed all five jobs in GitHub Actions run \href{https://github.com/miraflor/net-center/actions/runs/45778162775}{45778162775} (Python 3.10--3.13 and distribution build/smoke test). The Python 3.12 log reports 201 passing tests. The installed wheel was imported from site-packages and declared dependencies were checked. These are observations for that revision, not claims about future runs.

The original 111 tests cover analytical examples, dense edge sampling, pruning bounds, precision, parallel invariance, routing and topology. The 90 adaptive tests cover geometry and augmentation, component policies and group independence, exact and heuristic finite candidates, selection, connectors, audits and actual legacy integration. Numerical agreement with dense sampling is not an analytical proof of the continuous optimum; the envelope argument supplies the mathematical basis. Heuristic results require comparison against exact references on small cases, not claims of universal optimality. A correct solver cannot repair erroneous road connectivity or an inappropriate model.

Record the code revision, dependency versions, CRS, input data and provenance, topology and snapping options, modeled coverage, candidate universe, objectives and connectors, site-count policy, caps, seed and solver paths. Use \code{docs/ADAPTIVE_LOCATION.md} for API details, \code{docs/VALIDATION.md} for test instructions, and \code{docs/ALGORITHMS.md} for the original detailed algorithmic lineage and historical review. The earlier extensive single-center technical note is preserved in Git at \href{https://github.com/miraflor/net-center/tree/df24244e81e17b79e497b979fb3773382c87fed7/docs}{the previous source revision}. This unified note does not replace the historical audit record.

\begin{thebibliography}{99}
\bibitem{dijkstra} Dijkstra, E. W. (1959). A note on two problems in connexion with graphs. \emph{Numerische Mathematik}, 1, 269--271.
\bibitem{hakimi64} Hakimi, S. L. (1964). Optimum locations of switching centers and the absolute centers and medians of a graph. \emph{Operations Research}, 12(3), 450--459. \href{https://doi.org/10.1287/opre.12.3.450}{doi:10.1287/opre.12.3.450}.
\bibitem{hakimi65} Hakimi, S. L. (1965). Optimum distribution of switching centers in a communication network and some related graph theoretic problems. \emph{Operations Research}, 13(3), 462--475. \href{https://doi.org/10.1287/opre.13.3.462}{doi:10.1287/opre.13.3.462}.
\bibitem{kariv} Kariv, O., \& Hakimi, S. L. (1979). An algorithmic approach to network location problems. I: The p-centers. \emph{SIAM Journal on Applied Mathematics}, 37(3), 513--538. \href{https://doi.org/10.1137/0137040}{doi:10.1137/0137040}.
\bibitem{handler} Handler, G. Y., \& Mirchandani, P. B. (1979). \emph{Location on Networks: Theory and Algorithms}. MIT Press.
\bibitem{teitz} Teitz, M. B., \& Bart, P. (1968). Heuristic methods for estimating the generalized vertex median of a weighted graph. \emph{Operations Research}, 16(5), 955--961. \href{https://doi.org/10.1287/opre.16.5.955}{doi:10.1287/opre.16.5.955}.
\bibitem{kaufman} Kaufman, L., \& Rousseeuw, P. J. (1990). \emph{Finding Groups in Data}. Wiley.
\bibitem{gap} Tibshirani, R., Walther, G., \& Hastie, T. (2001). Estimating the number of clusters in a data set via the Gap statistic. \emph{JRSS B}, 63(2), 411--423. \href{https://doi.org/10.1111/1467-9868.00293}{doi:10.1111/1467-9868.00293}.
\bibitem{clusgap} R \texttt{cluster} package. \href{https://search.r-project.org/CRAN/refmans/cluster/html/clusGap.html}{\texttt{clusGap} documentation}, particularly \code{d.power} and \code{Tibs2001SEmax}.
\end{thebibliography}
\end{document}
